% Authors: Keerthi Rapolu & Sreeja Katta

The F1 improvement from 0.605 to 0.761 is driven primarily by recall. The
IFS-based detector captures anomalous workloads that the CPU-threshold baseline
classifies as normal because their utilization does not breach the detection
boundary. This pattern is consistent with the design intent of IFS: a workload
that finishes early and idles, consumes fewer resources than requested, or
executes a materially different computation than its description implies can
remain below the utilization alert threshold while still producing measurable
intent-behavior divergence.

The mismatch-subgroup result adds an important qualification. Among workloads
flagged as type-mismatched by intent inference, anomaly rates are higher and
mean IFS is lower than in the non-mismatch group, consistent with the hypothesis
that declared-to-inferred type divergence predicts runtime divergence. However,
the anomaly rate difference across subgroups is not large enough for type
mismatch alone to serve as a reliable detection signal. The subgroup result
supports IFS as a composite detector rather than suggesting that any single
structural feature is sufficient.

The precision tradeoff is also worth stating directly. IFS precision on this
benchmark is lower than the CPU-threshold baseline, meaning intent-aware
detection produces more false positives per true positive. This is a consequence
of IFS operating at broader sensitivity: it flags misaligned workloads earlier
and more broadly than a utilization gate would. The threshold sweep shows that an
operator can recover precision by raising $\theta_{\mathrm{ifs}}$ at the cost of
recall, which provides a practical tuning handle once deployed against a
production workload population.
